\documentclass[a4paper,11pt]{article}
\usepackage{ctex}
\usepackage{amsmath, amssymb}
\usepackage{booktabs}
\usepackage{float}
\usepackage{geometry}
\usepackage{array}
\usepackage{longtable}
\geometry{left=2.3cm, right=2.3cm, top=2.4cm, bottom=2.4cm}
\newcommand{\file}[1]{\texttt{\small #1}}

\begin{document}

\begin{center}
    \Large\textbf{上机作业 8：隐式 QR 算法与对称 QR 算法}
\end{center}

\section{题目与算法说明}
本次实验完成教材第 202 页第 2 题，以及第 244 页上机习题 1(2)、2(2)。
一般实矩阵的全部特征值和特征向量用隐式 QR 思想求解；实对称矩阵则采用对称 QR 方法求解。
程序文件为 \file{hw8.py}，结果数据写入以下文件：
\begin{itemize}
    \item \file{hw8\_p202\_roots.csv}；
    \item \file{hw8\_p202\_matrix\_eigenpairs.csv}；
    \item \file{hw8\_p244\_1\_summary.csv}；
    \item \file{hw8\_p244\_1\_eigenpairs.csv}；
    \item \file{hw8\_p244\_2\_eigenpairs.csv}。
\end{itemize}

对一般实矩阵 $A$，程序先用 Householder 相似变换化为上 Hessenberg 矩阵
\[
    A=Q_0HQ_0^T .
\]
随后在活动子块上进行带位移 QR 迭代
\[
    H_k-\mu_k I=Q_kR_k,\qquad
    H_{k+1}=R_kQ_k+\mu_k I,
\]
并用次对角元
\[
    |h_{i+1,i}| \le \varepsilon
    \bigl(|h_{ii}|+|h_{i+1,i+1}|+\|H\|_\infty\bigr)
\]
作为分裂判据。教材中隐式双步 QR 的核心思想是利用隐式 $Q$ 定理，在不显式形成
$M=H^2-sH+tI$ 的情况下完成同一个正交相似迭代。这里程序按等价的带位移 QR 形式实现，
矩阵阶数不大，便于直接检查每一步的残差；特征向量由
\[
    (A-\lambda I)v=0
\]
的最小奇异向量求得并归一化。

对实对称矩阵，程序先化为对称三对角矩阵
\[
    A=Q_0TQ_0^T,
\]
再对活动三对角子块使用 Wilkinson 位移
\[
    \mu=d_m-\frac{\operatorname{sign}(\delta)e_{m-1}^2}
    {|\delta|+\sqrt{\delta^2+e_{m-1}^2}},\qquad
    \delta=\frac{d_{m-1}-d_m}{2}.
\]
每一步累积正交变换，因此最终得到的特征向量列向量构成正交矩阵。所有数值结果均以后验相对残差
\[
    r_j=\frac{\|Av_j-\lambda_jv_j\|_\infty}
    {\max(\|A\|_\infty\|v_j\|_\infty,1)}
\]
检验；对称情形还检查 $\|V^TV-I\|_\infty$。

\section{教材第 202 页第 2 题}
\subsection{方程 $x^{41}+x^3+1=0$ 的全部根}
令
\[
    p(x)=x^{41}+x^3+1 .
\]
程序构造 Frobenius 伴随矩阵 $C$，使得 $\det(xI-C)=p(x)$，再对 $C$ 使用上面的
Hessenberg QR 过程。共进行 124 次 QR 迭代，全部根的最大缩放多项式残差为
$2.772\times 10^{-12}$。表 \ref{tab:p202-roots} 给出全部 41 个根；完整双精度数据见
\verb|hw8_p202_roots.csv|。

\begin{longtable}{rccc}
\caption{$x^{41}+x^3+1=0$ 的全部根}\label{tab:p202-roots}\\
\toprule
编号 & 根 & 模 & scaled residual \\
\midrule
\endfirsthead
\toprule
编号 & 根 & 模 & scaled residual \\
\midrule
\endhead
1 & $-0.968140341518-0.120866955421i$ & $0.975655954621$ & $4.224\times10^{-15}$ \\
2 & $-0.956339011328-0.273776209087i$ & $0.994755104159$ & $2.343\times10^{-14}$ \\
3 & $-0.910511134934-0.425528151316i$ & $1.005039668073$ & $3.969\times10^{-13}$ \\
4 & $-0.836863004991-0.567825669085i$ & $1.011318782380$ & $1.653\times10^{-14}$ \\
5 & $-0.739100565191-0.695904353165i$ & $1.015161324234$ & $9.163\times10^{-15}$ \\
6 & $-0.620672505471-0.805889307059i$ & $1.017198080159$ & $7.827\times10^{-15}$ \\
7 & $-0.485280239300-0.894538369840i$ & $1.017691409893$ & $1.347\times10^{-14}$ \\
8 & $-0.336984323231-0.959227612653i$ & $1.016698601838$ & $1.072\times10^{-14}$ \\
9 & $-0.180206043239-0.997962232726i$ & $1.014101984993$ & $1.029\times10^{-14}$ \\
10 & $-0.019728632208-1.009352156874i$ & $1.009544944772$ & $8.241\times10^{-15}$ \\
11 & $0.139165434976-0.992476733701i$ & $1.002186152982$ & $1.027\times10^{-13}$ \\
12 & $0.289812131068-0.946423674669i$ & $0.989802426390$ & $2.845\times10^{-13}$ \\
13 & $0.417151578256-0.871067309516i$ & $0.965802100303$ & $1.103\times10^{-13}$ \\
14 & $0.507568639468-0.810573438853i$ & $0.956376088965$ & $2.702\times10^{-13}$ \\
15 & $0.632339827064-0.753401199924i$ & $0.983599016337$ & $5.082\times10^{-13}$ \\
16 & $0.753719530443-0.655380276003i$ & $0.998807507353$ & $1.674\times10^{-12}$ \\
17 & $0.855158027644-0.532633535496i$ & $1.007468974896$ & $2.772\times10^{-12}$ \\
18 & $0.933664044732-0.392546384549i$ & $1.012828323284$ & $1.681\times10^{-12}$ \\
19 & $0.987183858032-0.240354383679i$ & $1.016022735628$ & $2.267\times10^{-14}$ \\
20 & $1.014304667340-0.080923029855i$ & $1.017527638416$ & $2.838\times10^{-14}$ \\
21 & $1.014304667340+0.080923029855i$ & $1.017527638416$ & $1.582\times10^{-14}$ \\
22 & $0.987183858032+0.240354383679i$ & $1.016022735628$ & $1.653\times10^{-14}$ \\
23 & $0.933664044732+0.392546384549i$ & $1.012828323284$ & $1.251\times10^{-14}$ \\
24 & $0.855158027644+0.532633535495i$ & $1.007468974896$ & $1.630\times10^{-13}$ \\
25 & $0.753719530443+0.655380276003i$ & $0.998807507353$ & $1.278\times10^{-12}$ \\
26 & $0.632339827064+0.753401199924i$ & $0.983599016337$ & $3.151\times10^{-15}$ \\
27 & $0.507568639468+0.810573438853i$ & $0.956376088965$ & $1.440\times10^{-15}$ \\
28 & $0.417151578256+0.871067309516i$ & $0.965802100303$ & $5.798\times10^{-16}$ \\
29 & $0.289812131068+0.946423674669i$ & $0.989802426390$ & $1.166\times10^{-12}$ \\
30 & $0.139165434976+0.992476733701i$ & $1.002186152982$ & $4.119\times10^{-13}$ \\
31 & $-0.019728632208+1.009352156874i$ & $1.009544944772$ & $6.302\times10^{-15}$ \\
32 & $-0.180206043239+0.997962232726i$ & $1.014101984993$ & $6.397\times10^{-15}$ \\
33 & $-0.336984323231+0.959227612653i$ & $1.016698601838$ & $1.525\times10^{-14}$ \\
34 & $-0.485280239300+0.894538369840i$ & $1.017691409893$ & $4.741\times10^{-15}$ \\
35 & $-0.620672505471+0.805889307059i$ & $1.017198080159$ & $1.922\times10^{-14}$ \\
36 & $-0.739100565191+0.695904353165i$ & $1.015161324234$ & $1.323\times10^{-14}$ \\
37 & $-0.836863004991+0.567825669085i$ & $1.011318782380$ & $3.364\times10^{-14}$ \\
38 & $-0.910511134934+0.425528151316i$ & $1.005039668073$ & $1.247\times10^{-14}$ \\
39 & $-0.956339011328+0.273776209087i$ & $0.994755104159$ & $2.271\times10^{-13}$ \\
40 & $-0.968140341518+0.120866955421i$ & $0.975655954621$ & $7.388\times10^{-14}$ \\
41 & $-0.952483875214$ & $0.952483875214$ & $9.992\times10^{-15}$ \\
\bottomrule
\end{longtable}

\subsection{参数矩阵的全部特征值和特征向量}
教材第 202 页第 2(3) 题中的矩阵为
\[
A(x)=
\begin{bmatrix}
9.1&3.0&2.6&4.0\\
4.2&5.3&4.7&1.6\\
3.2&1.7&9.4&x\\
6.1&4.9&3.5&6.2
\end{bmatrix}.
\]
表 \ref{tab:p202-param} 给出 $x=0.9,1.0,1.1$ 时的全部特征值。
每个特征值对应的特征向量列在随附 CSV 文件中。

\begin{table}[H]
\centering
\small
\caption{参数矩阵 $A(x)$ 的全部特征值}
\label{tab:p202-param}
\begin{tabular}{clc}
\toprule
$x$ & 特征值 & 最大相对残差 \\
\midrule
0.9 & $2.870401735770\pm0.642891129472i,\ 6.819518337521,\ 17.439678190939$
& $3.862\times10^{-16}$ \\
1.0 & $2.867999246456\pm0.688747355259i,\ 6.787516591559,\ 17.476484915529$
& $1.646\times10^{-16}$ \\
1.1 & $2.865457870674\pm0.732169739115i,\ 6.756056187178,\ 17.513028071474$
& $4.291\times10^{-16}$ \\
\bottomrule
\end{tabular}
\end{table}

由表可见，当参数 $x$ 从 $0.9$ 增至 $1.1$ 时，两个实特征值分别约从
$6.8195$ 降至 $6.7561$、从 $17.4397$ 升至 $17.5130$；共轭复特征值的实部略降，
虚部绝对值从 $0.6429$ 增至 $0.7322$。这说明该参数主要扩大了这对复根的虚部。

\section{教材第 244 页上机习题 1(2)}
本题矩阵为 $n=50,51,\ldots,100$ 阶实对称三对角矩阵
\[
    A_n=\operatorname{tridiag}(1,4,1).
\]
程序逐个阶数调用对称 QR 子程序计算全部特征值和正交特征向量。
完整数据保存在随附 CSV 文件中。
为了核验结果，同时使用该 Toeplitz 三对角矩阵的解析表达式
\[
    \lambda_k=4+2\cos\frac{k\pi}{n+1},\qquad
    q_j^{(k)}=\sqrt{\frac{2}{n+1}}\sin\frac{jk\pi}{n+1},
    \quad 1\le j,k\le n .
\]
程序输出的 QR 特征值与解析值的最大差异约为 $10^{-15}$ 量级。

\begin{table}[H]
\centering
\small
\caption{$A_n=\operatorname{tridiag}(1,4,1)$ 的对称 QR 结果摘要}
\label{tab:p244-1}
\begin{tabular}{rrrrrr}
\toprule
$n$ & $\lambda_{\min}$ & $\lambda_{\max}$ & QR 迭代 & 最大残差 & $\|V^TV-I\|_\infty$ \\
\midrule
50  & 2.003793342526 & 5.996206657474 & 101 & $3.925\times10^{-14}$ & $1.727\times10^{-14}$ \\
60  & 2.002651820230 & 5.997348179770 & 123 & $3.942\times10^{-14}$ & $2.039\times10^{-14}$ \\
70  & 2.001957546960 & 5.998042453040 & 142 & $5.228\times10^{-14}$ & $2.858\times10^{-14}$ \\
80  & 2.001504094992 & 5.998495905008 & 159 & $4.909\times10^{-14}$ & $2.794\times10^{-14}$ \\
90  & 2.001191718898 & 5.998808281102 & 181 & $2.266\times10^{-14}$ & $3.145\times10^{-14}$ \\
100 & 2.000967435416 & 5.999032564584 & 200 & $2.185\times10^{-14}$ & $3.481\times10^{-14}$ \\
\bottomrule
\end{tabular}
\end{table}

从 $n=50$ 到 $n=100$，最小特征值逐渐趋近于 $2$，最大特征值逐渐趋近于 $6$。
这是因为 $\cos(k\pi/(n+1))$ 的取值随 $n$ 增大更接近区间 $[-1,1]$ 的端点。
对称 QR 产生的特征向量保持良好正交性，表中最大正交性误差均在 $10^{-14}$ 量级。

\section{教材第 244 页上机习题 2(2)}
本题矩阵为 $100$ 阶
\[
    B=\operatorname{tridiag}(-1,2,-1).
\]
对称 QR 子程序给出的最小、最大特征值及对应特征向量残差见表 \ref{tab:p244-2}。

\begin{table}[H]
\centering
\caption{$B=\operatorname{tridiag}(-1,2,-1)$ 的最大和最小特征值}
\label{tab:p244-2}
\begin{tabular}{lcc}
\toprule
特征值 & 数值 & 相对残差 \\
\midrule
$\lambda_{\min}$ & 0.0009674354160239 & $2.802\times10^{-16}$ \\
$\lambda_{\max}$ & 3.9990325645839762 & $4.996\times10^{-16}$ \\
\bottomrule
\end{tabular}
\end{table}

该矩阵同样有解析校验式
\[
    \lambda_k=2-2\cos\frac{k\pi}{101},\qquad
    q_j^{(k)}=\sqrt{\frac{2}{101}}\sin\frac{jk\pi}{101},\quad 1\le j,k\le 100.
\]
因此
\[
    \lambda_{\min}=\lambda_1=2-2\cos\frac{\pi}{101},\qquad
    \lambda_{\max}=\lambda_{100}=2-2\cos\frac{100\pi}{101}.
\]
对应单位特征向量可写为
\[
    q_j^{(1)}=\sqrt{\frac{2}{101}}\sin\frac{j\pi}{101},
\]
\[
    q_j^{(100)}=\sqrt{\frac{2}{101}}\sin\frac{100j\pi}{101}
    =(-1)^{j+1}\sqrt{\frac{2}{101}}\sin\frac{j\pi}{101}.
\]
程序输出的完整 100 个分量保存在 \file{hw8\_p244\_2\_eigenpairs.csv}。

\section{结论}
本次实验实现了两个通用特征值计算子程序。对于一般实矩阵，先作 Hessenberg 约化再用带位移
QR 迭代分裂出全部特征值，并用最小奇异向量求特征向量；对于实对称矩阵，先三对角化再用
Wilkinson 位移对称 QR 迭代，同时累积正交变换得到正交特征向量。

数值结果显示：$x^{41}+x^3+1=0$ 的全部根均被求出，最大缩放残差为
$2.772\times10^{-12}$；第 202 页参数矩阵的全部特征对相对残差在 $10^{-16}$ 量级；
第 244 页两个实对称三对角问题的特征向量正交性和特征残差均达到双精度舍入误差附近。
因此程序满足本次上机作业对全部特征值和特征向量计算的精度要求。

\end{document}
